\documentclass{article}
\usepackage[utf8]{inputenc}
\usepackage[T1]{fontenc}
\usepackage{geometry}
\geometry{
    a4paper,
    left=2.5cm,
    right=2.5cm,
    top=2.5cm,
    bottom=2.5cm
}
\usepackage{graphicx}
\usepackage{float}
\usepackage{hyperref}
\usepackage[english]{babel}
\usepackage{amsmath}
\usepackage{amssymb}
\usepackage{array}
\usepackage{booktabs}
\usepackage{caption}
\usepackage{subcaption}
\usepackage{enumitem}
\usepackage{fancyhdr}
\usepackage{titlesec}
\usepackage{url}
\usepackage{textcomp}
\usepackage{xcolor}
\usepackage{listings}
\usepackage{multirow}
\usepackage{longtable}
\usepackage{tabularx}
\usepackage{times}

% Set up page style
\pagestyle{fancy}
\fancyhf{}
\rfoot{\thepage}
\renewcommand{\headrulewidth}{0pt}

% Title formatting
\titleformat{\chapter}[display]{\normalfont\bfseries}{\chaptertitle\thechapter}{0pt}{\Huge}
\titlespacing*{\chapter}{0pt}{50pt}{40pt}

% Document info
\title{Template Document}
\author{Academic Support Agent}
\date{\today}

\begin{document}
\maketitle

\maketitle


\begin{abstract}The abstract should provide clear information about the research and the results obtained, and should not exceed 200 words. The abstract should not contain citations.

\keywords{Please provide key words or phrases to enable retrieval and indexing. Acronyms should be avoided}
\end{abstract}


\section{Introduction}

It has been discovered that a number of elements (Nb, Ca, Pb, etc.) and 
compounds (NbN, MgB$_{2}$, FeSe, etc.) have exhibited superconducting 
behavior since the discovery of superconductivity \cite{1}.

\textcolor{red}{\textbf{Symbols, units, and abbreviations:}} In general, the journal follows the conventions of Scientific Style and Format, The
CSE Manual for Authors, Editors, and Publishers, Council of Science Editors, Reston, VA, USA (7th ed.). All abbreviations
and acronyms should be defined at first mention.

\section{Section}
Calculations are carried out by using the PWR TMI-1 unit cell case 
introduced in the UAM test problems. (\ref{eq1}).
\begin{equation}
\label{eq1}
equation
\end{equation}
Here, $T$ is the temperature of the fuel region (with superscript $F)$ and $\rho $ 
is the reactivity of the system.

\subsection{Subsection}

\textcolor{red}{\textbf{Manuscript content:}} Research articles should be divided into logically ordered and numbered sections. Principal
sections should be numbered consecutively with Arabic numerals (1. Introduction, 2. Formulation of problem, etc.) and
subsections should be numbered 1.1., 1.2., etc. Do not number the Acknowledgements or References sections.

It has been discovered that a number of elements (Nb, Ca, Pb, etc.) and 
compounds (NbN, MgB$_{2}$, FeSe, etc.) have exhibited superconducting 
behavior since the discovery of superconductivity \cite{2}. 
\begin{equation}
\label{eq2}
equation
\end{equation}

It has been discovered that a number of elements (Nb, Ca, Pb, etc.) and 
compounds (NbN, MgB$_{2}$, FeSe, etc.) have exhibited superconducting 
behavior since the discovery of superconductivity.

\subsubsection{Subsubsection}
It has been discovered that a number of elements (Nb, Ca, Pb, etc.) and 
compounds (NbN, MgB$_{2}$, FeSe, etc.) have exhibited superconducting 
behavior since the discovery of superconductivity \cite{1}. Table \ref{tab1}, Figure \ref{fig1}.

\textcolor{red}{\textbf{Tables and Figures}}

All illustrations (photographs, drawings, graphs, etc.), not including tables, must be labelled ``Figure."
All tables and figures must have a caption and/or legend and be numbered (e.g., Table 1, Figure 2), unless there is
only one table or figure, in which case it should be labelled ``Table" or ``Figure" with no numbering. Captions must be
written in sentence case (e.g., Jackknife and bootstrap with moving blocks.).
All tables and figures must be numbered consecutively as they are referred to in the text. Please refer to tables and
figures with capitalisation and unabbreviated (e.g., ``As shown in Figure 2 ", and not ``Fig. 2" or ``figure 2").
The resolution of images should not be less than 118 pixels/cm when width is set to 16 cm. Images must be scanned
at 1200 dpi resolution and submitted in jpeg or tiff format.
Graphs and diagrams must be drawn with a line weight between 0.5 and 1 point. Graphs and diagrams with a line
weight of less than 0.5 point or more than 1 point are not accepted. Scanned or photocopied graphs and diagrams are not
accepted.
Charts must be prepared in 2 dimensions unless required by the data used. Charts unnecessarily prepared in 3
dimensions are not accepted.

\begin{table}[h!]
\caption{default.}
\medskip
\centering\renewcommand{\arraystretch}{1.2}
\begin{tabular}{|l|l|l|l|}
\hline
Parameters& 
S0& 
S50& 
S100 \\
\hline
a ({\AA})& 
3819& 
3811& 
3796 \\
\hline
b ({\AA})& 
3890& 
3875& 
4013 \\
\hline
c ({\AA})& 
11,651& 
11,648& 
11,617 \\
\hline
T$_{c}$ (K)& 
85& 
88& 
88 \\
\hline
J$_{c}$ (0) (A/cm$^{2})$& 
9489& 
14,287& 
12,817 \\
\hline
J$_{c,\max}$ (A/cm$^{2})$& 
12,571& 
17,751& 
15,156 \\
\hline
\end{tabular}
\label{tab1}
\centering
\end{table}

\begin{figure}[h!]
\centering
\includegraphics[width=17cm]{figure_example.eps}\vs{-2mm}
\caption{default.}
\label{fig1}
\end{figure}


\section{Conclusion}
It has been discovered that a number of elements (Nb, Ca, Pb, etc.) and 
compounds (NbN, MgB$_{2}$, FeSe, etc.) have exhibited superconducting 
behavior since the discovery of superconductivity.

\section*{Acknowledgment}
Acknowledgement and/or disclaimers, if any\\
Names of funding organizations should be written in full.

\textbf{References}

Citations in the text should be identified by numbers in square brackets. If the manuscripts are prepared in LaTeX format, the citations should be given via cite command during the preparation of LaTeX file, not given in square brackets manually. The list of references at the end of the
paper should be given in order of their first appearance in the text. All authors should be included in reference lists
unless there are 6 or more, in which case only the first 5 should be given, followed by et al.. Do not use individual
sets of square brackets for citation numbers that appear together, e.g., \cite{2,3,4,5,6}, not \cite{2}. Do not include
personal communications, unpublished data, websites, or other unpublished materials as references, although such
material may be inserted (in parentheses) in the text. References should be formatted as follows (please note the
punctuation and capitalisation):

\begin{thebibliography}{99}
\bibitem{1} Saleh ZM, Kmail SM, Assaf SF, Qasrawi AF. Recombination mechanisms in hydrogenated silicon nanocrystalline thin films. Turkish Journal of Physics 2013; 37 (3): 283-288. doi: 10.3906/ z-1301-12

\bibitem{2} Deschler F, Price M, Pathak S, Klintberg LE, Jarausch DD et al. High photoluminescence efficiency and optically pumped lasing in solution-processed mixed halide perovskite semiconductors. Journal of Physical Chemistry Letters 2014; 5 (8): 1421-1426. doi: 10.1021/jz5005285

\bibitem{3} Altintas A. Tristor ve triyak harmoniklerinin 3 boyutlu g sterimi ve toplam harmonik bozunuma egri uydurma. Pamukkale Universitesi Muhendislik Bilimleri Dergisi 2004; 10 (3): 415-421 (in Turkish).

\bibitem{4} Kocyigit Y, Dilmac S. Cardiac arrhythmia diagnosis using firefly algorithm. KS  Muhendislik Bilimleri Dergisi 2018; 21 
(3): 226-234 (in Turkish with an abstract in English). doi: 10.17780/ksujes.435734

\bibitem{5} Hedenmalm H, Korenblum B, Zhu K. Theory of Bergman Spaces. New York, NY, USA: Springer-Verlag, 2000.

\bibitem{6} Lui SH. Numerical Analysis of Partial Differential Equations. Hoboken, NJ, USA: Wiley, 2011.

\bibitem{7} Gries P, Campbell J, Montojo J. Functions that Python provides. In: Coron T (editor). Practical Programming: An Introduction to Computer Science Using Python 3.6. 3rd ed. Princeton, NJ, USA: Princeton University Press, 2012, pp. 31-33.

\bibitem{8} Luszczek P, Dongarra J. Design of interactive environment for numerically intensive parallel linear algebra calculations. In: Computational Science  4th International Conference Proceedings; Krakow, Poland; 2004. pp. 270-277.

\bibitem{9} McKenna FT. Object-oriented finite element programming: frameworks for analysis, algorithms and parallel computing. PhD, University of California, Berkeley, CA, USA, 1997.
\end{thebibliography}
\end{document}